\documentclass[11pt]{article}
\usepackage[margin=1in]{geometry}
\usepackage[T1]{fontenc}
\usepackage{lmodern}
\usepackage{booktabs}
\usepackage{enumitem}
\usepackage{graphicx}
\usepackage{hyperref}
\usepackage{xcolor}
\usepackage{amssymb}
\newcommand{\status}[1]{\textcolor{red}{\textbf{#1}}}
\title{Subject-Held-Out Probability Calibration on HARTH: An Exploratory Custody-Released Analysis}
\author{HowHow Research}
\date{Exploratory report v2.1 --- 2026-08-24}
\begin{document}
\maketitle
\begin{center}\status{EXPLORATORY / POST-OBSERVATION / INCONCLUSIVE --- NOT A CONFIRMATORY CLAIM}\end{center}
\begin{abstract}
We report a custody-released, exploratory analysis of a deterministic CPU nearest-centroid activity-recognition baseline on HARTH. The frozen population contains twenty-two subjects and sixty-six configuration-folds. Nested subject-held-out evaluation compares uncalibrated probabilities with inner-subject temperature scaling for full, back-only, and thigh-only sensor representations. Subject-macro uncertainty uses 2,000 bootstrap replicates; paired sign-flip tests use 200,000 draws. Full-sensor calibrated NLL is generated in Table~\ref{tab:metrics}; the unadjusted paired NLL p-value is generated in the evidence snapshot, but Holm adjustment yields no familywise rejection. Brier and ECE do not establish a consistent calibration benefit. Exploratory support-aware F1 preserves \texttt{NOT\_ESTIMABLE}. This is an inconclusive post-observation report: it makes no novelty, causal, superiority, transferability, or generalization claim.
\end{abstract}

\section{Introduction}
Probability quality matters when activity-recognition systems are used to communicate uncertainty, but a held-out-subject estimate is not a deployment guarantee. This report answers a narrow question: how did a frozen deterministic CPU baseline behave under the released HARTH run3 analysis, with and without inner-subject temperature scaling? The analysis is deliberately bounded to a frozen population and declared estimands. It does not estimate clinical utility, real-time safety, or population-wide calibration.

\section{Research Questions and Claim Boundary}
\begin{enumerate}[label=RQ\arabic*:]
\item What are subject-macro NLL, multiclass Brier score, and ten-bin top-label ECE for the uncalibrated baseline?
\item What paired changes accompany inner-subject temperature scaling without using outer-test labels?
\item What descriptive differences appear for back-only and thigh-only sensor ablations?
\end{enumerate}
All results are exploratory and post-observation. The result's custody boundary is \texttt{guarded\_real\_quarantined\_no\_release}; custody acceptance authorizes manuscript generation but does not release \texttt{quarantine.json}. The report preserves negative and inconclusive evidence and does not convert an artifact gate into scientific or reviewer acceptance.

\section{Methods}
\subsection{Design and estimands}
The primary estimand is the subject-macro expected test loss under leave-one-subject-out evaluation. All windows from a held-out subject remain in that subject's outer fold. The frozen population and three sensor configurations yield sixty-six lifecycle folds. The baseline is deterministic nearest-centroid classification over channel means and standard deviations. Calibration selects a bounded scalar temperature only with inner training-subject validation, then applies it once to the outer test subject.

\subsection{Metrics and uncertainty}
NLL, multiclass Brier score, and ten-bin top-label ECE are summarized per subject and then macro-averaged. The generator reads the strict result and emits every displayed numeric cell; no metric is hand-entered in this manuscript. Intervals are percentile subject-cluster bootstrap intervals using 2,000 replicates. The three calibration comparisons use Holm adjustment at the declared familywise threshold. Sign-flip tests use 200,000 PCG64 draws with inclusive ties and unchanged zero handling. Sensor ablations are descriptive.

\subsection{Custody and reproducibility contract}
The builder verifies the exact released result SHA, schema and COMPLETE state, frozen population, exploratory boundary, and custody flags before deriving a self-contained snapshot. It strips source destination and private-path fields, retains result JSON pointers, and binds each derived cell to the result SHA. `--check` compares snapshot, tables, and figures byte-for-byte. Failure on any release, hash, schema, status, value, pointer, or reproducibility drift is intentional.

\section{Results}
The generated tables report all full, reduced-sensor, calibrated, and uncalibrated estimates with intervals. Figure~\ref{fig:nll-bars} visualizes the full-sensor NLL estimates and Figure~\ref{fig:nll-delta} visualizes their paired contrast. The full-sensor contrast is compatible with a pre-adjustment NLL difference, while the Holm-adjusted family decision is a non-rejection. Brier and ECE contrasts do not provide a consistent calibration pattern. Reduced-sensor NLL increases are descriptive contrasts, not sensor-attribution effects.

% Generated by tools/generate_tables.py. Do not edit.
\section{Results Placeholder}
\begin{table}[ht]
\centering
\caption{Protocol-v2 metrics.}
\label{tab:results}
\fbox{\parbox{0.82\linewidth}{\centering\Large UNVERIFIED\\[4pt]\normalsize No validator-approved v2 result artifact was supplied.\\Metrics are intentionally withheld.}}
\end{table}
\begin{figure}[ht]
\centering
\fbox{\parbox{0.72\linewidth}{\centering\Large UNVERIFIED\\[4pt]\normalsize Figure generation is closed until validated fold-level artifacts exist.}}
\caption{Protocol-v2 result figure placeholder.}
\label{fig:results}
\end{figure}

% GENERATED FIGURES: dimensions are derived from evidence-snapshot.json.
\begin{figure}[ht]
\centering
\fbox{\begin{minipage}{.86\linewidth}\small
Full-sensor NLL (lower is better)\\[2pt]
Calibrated\quad\rule[2pt]{68pt}{7pt}\quad 0.9724\\
Uncalibrated\quad\rule[2pt]{80pt}{7pt}\quad 1.1402
\end{minipage}}\caption{Generated full-sensor calibration comparison.}\label{fig:nll-bars}
\end{figure}
\begin{figure}[ht]
\centering
\fbox{\begin{minipage}{.86\linewidth}\small
Calibration delta, full sensor (calibrated minus uncalibrated)\\[2pt]
\rule[2pt]{23pt}{7pt}\quad -0.1678\quad (negative favors calibrated NLL)
\end{minipage}}\caption{Generated paired NLL contrast; interval and multiplicity adjustment remain in the tables.}\label{fig:nll-delta}
\end{figure}


\section{Exploratory F1 and Support}
Classwise F1 is computed from TP, FP, and FN only when its denominator and support permit estimation. Zero support or zero denominator remains \texttt{NOT\_ESTIMABLE}; it is not a zero, an imputation, or evidence of failure. F1 is secondary exploratory context and is not used to support a calibration claim.

\section{Discussion}
The released evidence is balanced and inconclusive. Temperature scaling coincided with a negative full-sensor NLL contrast, but multiplicity adjustment removed familywise support, and Brier/ECE did not show a consistent benefit. Ablations show larger descriptive NLL differences for reduced representations, but the design cannot attribute those differences causally to a sensor. These observations do not establish novelty, superiority, transferability, or generalization.

\section{Limitations and Lifecycle Disclosure}
Only the frozen twenty-two subjects and sixty-six folds are represented. The custody review found that `final.json` says COMPLETE while reporting `scientific\_metrics=false` and `completed\_folds=[]`; the completed engine snapshot records the lifecycle folds. This minor, disclosable metadata defect is not silently reconciled and is not treated as scientific evidence. The retained package lacks subject-level sufficient statistics and sign counts, so exact independent regeneration of every bootstrap interval and paired-test value is unavailable. The initial custody dissent and the limitation are retained in the external adjudication record. `quarantine.json` remains unchanged and is not claimed as released or mutated.

Stronger evidence would require prospective hypotheses, corrected lifecycle metadata, governed retention of sufficient statistics and sign counts, and independent validation. None is performed here; run3 is not rerun or tuned.

\section{Provenance and Availability}
The tracked source manifest excludes raw/private artifacts. The generated snapshot records the immutable result hash and source JSON pointers. Code, protocol, schema, and publication builder remain in the repository; raw data, private custody paths, and the quarantined package are not copied into this paper. MiKTeX compilation is a manuscript-mechanics gate only, not scientific validation.

\bibliographystyle{plain}
\bibliography{references}
\end{document}
